% Fragmento compartido por VII y CRISTAL; incluir despues del cuerpo Feigenbaum.
% Procedencia: resultados recuperados, con demostraciones elementales reunidas.
% Perfil, Jacobi y escalado: 00_antecedentes_operatorios.tex, 03, 06 y 07
% del presente paquete; no se cambia su selector ni su dominio uniforme.
% Constitutiva: CRISTAL_ES/main_autosuficiente.tex:58738--58840,58864--58894.
% Inversion: CRISTAL_ES/sections/cr28_restitucion_constitutiva.tex:18--59.
% Catalan: el mismo propietario, 277--373; identidad integral, unicidad y serie.
% Corte material: CUMPLIMIENTO_EDITORIAL_20260928/fuentes/CRISTAL_ES.
% No contiene una nueva calibracion ni una nueva campana de calculo.
\section{Articulation of the critical, spectral and constitutive readers}
\label{hmtfg:articulacion:seccion}

The APP--TRIT--TPK composition precedes the readers in this section.
APP preserves the additive and multiplicative evaluations in each cell,
with their remainders and quotients; tritic reduction determines the regime,
orientation and carry of the transition. The effective composition
of selection, transport and memory update produces the enriched
state and its prolongations, developed in
\ref{hmtfg:app-trit-transporte} and
\ref{hmtfg:subsubsec:tres-vueltas-ext-ch}.
Nonadic holonomy composes those transitions: phase return
advances memory and preserves history. The readers act on the
same joint discrete structure of the continuum, with its five
consubstantial constructions and with the two readouts of the same
antecedent: ordered evaluation and boundary incidence.

The oriented wheel and angular channels are outputs of this
composition. Their already generated HMT coordinates are used in subsequent
operations, with their marks and domains; they are not metrological numbers
introduced to select an answer. The ordered readout
allows the following formulas to be evaluated, while the provenance state
preserves sheet, route, remainder, quotient, carry, boundary and memory.
The recoveries that are proved recover the object of their declared
domain; the evaluated representation remains accompanied by those
records.

\subsection{Critical separation as an intrinsic asymptotic reading}
\label{hmtfg:articulacion:separacion}

In the uniform sector of the dodecaphase wheel, the oriented
representatives and their weights are
\[
 \phi_m=\frac{(3m-1)\pi_{\mathrm{HMT}}}{18},\qquad
 w_m=\frac1{12},\qquad 1\le m\le12.
\]
The second moment is normalized through
\(\sum_{m=1}^{12}(3m-1)^2=5394=6\cdot899\).
Thus the common angular factor cancels, leaving
\begin{equation}
 k_m=\frac{2(3m-1)}{\sqrt{899}},\qquad
 u_0(z)=\frac1{12}\sum_{m=1}^{12}\bigl(1-\cos(k_mz)\bigr),
 \qquad f_\lambda(z)=1-\lambda u_0(z).
 \label{hmtfg:articulacion:perfil}
\end{equation}
Variable \(\lambda\) ranges over the family after the
reader, its weights and normalization have been fixed. In
\(0<\lambda\le2/u_0(1)\), the dynamical domain is \([-1,1]\).
The first root is \(\lambda_1=1/u_0(1)\); for \(n\ge2\),
\(\lambda_n\) is the smallest root greater than \(\lambda_{n-1}\)
whose critical point has exact period \(2^n\).
The classification, existence of each minimum and metric
identification of these centers have been proved in
\ref{hmtfg:escalado:15}--\ref{hmtfg:escalado:17}.

\begin{proposition}[Structural reading of the scaling constant]
\label{hmtfg:articulacion:delta}
For the uniform reader \eqref{hmtfg:articulacion:perfil} and its
first-root selector, the critical-separation constant is
\begin{equation}
 \delta_{\mathrm{HMT,crit}}
 :=\lim_{n\to\infty}
 \frac{\lambda_n-\lambda_{n-1}}{\lambda_{n+1}-\lambda_n}
 =\delta_{\mathrm F}.
 \label{hmtfg:articulacion:limite}
\end{equation}
It is the asymptotic parameter-separation invariant of the critical
returns of that reader, with its selection order preserved.
\end{proposition}
\begin{proof}
The scaling theorem \eqref{hmtfg:escalado:eq:67} provides
\(\lambda_*-\lambda_n=C\delta_{\mathrm F}^{-n}(1+e_n)\),
with \(C>0\) and \(e_n\to0\). Subtracting two consecutive terms,
\[
 \lambda_n-\lambda_{n-1}
 =C\delta_{\mathrm F}^{-n}
   \bigl[\delta_{\mathrm F}-1+
          \delta_{\mathrm F}e_{n-1}-e_n\bigr].
\]
The quotient of this equality and the one at index \(n+1\) converges to
\(\delta_{\mathrm F}\). The eventual identity of the local centers
with the first roots, proved in
\eqref{hmtfg:escalado:eq:66}, allows the scaling to be applied precisely
to the selected sequence and not to another cascade chosen afterward.
\end{proof}

This formulation preserves more information than the terminal numeral:
it identifies the wheel, its reading, the family and the selector whose limit
is evaluated. Its interpretation is intrinsic to that HMT composition;
recognition of the universal constant subsequently uses the analytic
theorem applied to the already constructed reader.

\subsection{The finite measure and spectral recovery of the profile}
\label{hmtfg:articulacion:jacobi}

The squared frequencies produce the atomic probability measure
\begin{equation}
 s_m=k_m^2=\frac{4(3m-1)^2}{899},\qquad
 \nu_0=\frac1{12}\sum_{m=1}^{12}\delta_{s_m},\qquad
 M_j=\int s^j\,d\nu_0(s).
 \label{hmtfg:articulacion:medida-finita}
\end{equation}
Here \(\delta_{s_m}\) is the Dirac mass at a node, distinct
from the constant \(\delta_{\mathrm F}\). The measure preserves the
twelve frequencies and their weights. For a nonzero polynomial of degree less
than \(r\le12\),
\[
 \int p(s)^2\,d\nu_0(s)
 =\frac1{12}\sum_{m=1}^{12}p(s_m)^2>0,
\]
since a polynomial of degree less than twelve cannot vanish at the
twelve distinct nodes. Thus the Hankel matrices
\((M_{i+j})_{0\le i,j<r}\) are positive definite.
The polynomial \(\prod_m(s-s_m)\) has zero norm; the space
\(L^2(\nu_0)\) has dimension twelve.

\begin{proposition}[Exact representation of the critical reader]
\label{hmtfg:articulacion:representacion}
Let \(\psi_0,\ldots,\psi_{11}\) be the orthonormal polynomials
with positive leading coefficient for \(\nu_0\), with
\(\psi_0=1\). Multiplication by \(s\) has Jacobi
matrix \(J_{12}\), real symmetric with positive subdiagonal entries.
With \(e_0=(1,0,\ldots,0)^{\mathsf T}\),
\begin{equation}
 u_0(z)=\int\bigl(1-\cos(z\sqrt{s})\bigr)\,d\nu_0(s)
 =e_0^{\mathsf T}\bigl[I-\cos(z\sqrt{J_{12}})\bigr]e_0.
 \label{hmtfg:articulacion:perfil-jacobi}
\end{equation}
\end{proposition}
\begin{proof}
The preceding positivity allows orthogonalization through degree eleven.
If \(i<j-1\), self-adjointness of multiplication gives
\(\langle s\psi_j,\psi_i\rangle
=\langle\psi_j,s\psi_i\rangle=0\).
By symmetry, only the diagonal and its two neighbors remain; the positive
ratio of the leading coefficients establishes positivity of the
subdiagonal. Define
\[
 Q_{mj}=\frac{\psi_j(s_m)}{\sqrt{12}},\qquad
 D=\operatorname{diag}(s_1,\ldots,s_{12}).
\]
Orthonormality gives \(Q^{\mathsf T}Q=I\), and evaluating the
recurrence at the nodes gives \(DQ=QJ_{12}\). Since \(Q\) is
square, \(J_{12}=Q^{\mathsf T}DQ\). Its eigenvalues are
positive, so the positive square root and cosine are calculated in
this diagonalization. Finally,
\(Qe_0=(1,\ldots,1)^{\mathsf T}/\sqrt{12}\), which proves
\eqref{hmtfg:articulacion:perfil-jacobi}.
\end{proof}

The realization through \((J_{12},e_0)\) recovers the same profile;
therefore, preserving the family and selector also recovers
its returns and first roots. Spectral positivity
justifies this representation. Parameter transversality
also uses the derivative of the return and its orbital products,
developed in \ref{hmtfg:escalado:14.3} and
\ref{hmtfg:escalado:6}. These are properties of different operations:
the scaling proof preserves both and does not replace
tangent control with positivity of the moments.

\subsection{Constitutive channels and their labeled inversion}
\label{hmtfg:articulacion:constitutiva}

The constitutive reader acts on another typed object on the same
support: two-sheet angular transport. Its already generated
coordinates are \(x=A_{\rm rad}\) and \(y=C^*_{\rm rad}\), in the
contracting chamber \(x>|y|\). The sheet involution
\(R=R^*=R^{-1}\) determines the marked projectors
\(P_\pm=(I\pm R)/2\). The transport representation is
\begin{equation}
 q_+=\exp[-(x+y)],\qquad q_-=\exp[-(x-y)],\qquad
 T=q_+P_++q_-P_-,\qquad 0<q_\pm<1.
 \label{hmtfg:articulacion:canales}
\end{equation}
Exponential evaluation realizes those channels after their
angular construction. The transport incidences \(90\) and \(120\)
fix the response
\[
 \mathcal V=(I-T^{90})(I-T^{120})^{-1}.
\]
The inverse exists because \(1-q_\pm^{120}>0\).
Writing \(s=q^{30}\), with \(30=\gcd(90,120)\), gives
\begin{equation}
 f(s)=\frac{1-s^3}{1-s^4}
     =\frac{1+s+s^2}{1+s+s^2+s^3},\qquad
 r_\pm=f(q_\pm^{30}),\qquad
 \mathcal V=r_+P_++r_-P_-.
 \label{hmtfg:articulacion:respuesta}
\end{equation}
In this subsection \(f\) denotes the constitutive response;
\(f_\lambda\) still denotes the dynamical map of
\eqref{hmtfg:articulacion:perfil}. Their arguments and domains
are distinct.

\begin{proposition}[Recovery of the channels and orientation]
\label{hmtfg:articulacion:inversion}
The constitutive function \(f:(0,1)\to(3/4,1)\) is a decreasing
bijection. Its inverse \(\psi\) assigns to \(r\) the unique root
\(s\in(0,1)\) of
\begin{equation}
 r s^3+(r-1)(1+s+s^2)=0.
 \label{hmtfg:articulacion:cubica}
\end{equation}
The normalized constitutive readings, with their labeled sheets,
\[
 \widehat\varepsilon=r_+^2,\qquad
 \widehat\mu=r_-^2,\qquad
 \widehat Z=\frac{r_-}{r_+},\qquad
 \widehat c=\frac1{r_+r_-},
\]
with \(\widehat\varepsilon,\widehat\mu\in(9/16,1)\), allow recovery of
\begin{equation}
 \begin{aligned}
 r_+&=\sqrt{\widehat\varepsilon},&
 r_-&=\sqrt{\widehat\mu},\\
 s_\pm&=\psi(r_\pm),& q_\pm&=s_\pm^{1/30}.
 \end{aligned}
 \label{hmtfg:articulacion:recuperacion-canales}
\end{equation}
\begin{equation}
 x=-\frac12\log(q_+q_-),\qquad
 y=\frac12\log\frac{q_-}{q_+}.
 \label{hmtfg:articulacion:recuperacion-angular}
\end{equation}
\end{proposition}
\begin{proof}
Differentiating the rational quotient gives
\begin{equation}
 f'(s)=-\frac{s^2(3+2s+s^2)}{(1+s+s^2+s^3)^2}<0
 \qquad(0<s<1).
 \label{hmtfg:articulacion:monotonia}
\end{equation}
The limits \(f(0^+)=1\) and \(f(1^-)=3/4\), together with
continuity, prove bijectivity. Multiplying \(f(s)=r\) by
the positive denominator produces \eqref{hmtfg:articulacion:cubica}.
Positivity of \(r_\pm\) determines the square roots
and that of \(s_\pm\) determines the thirtieth roots.
By \eqref{hmtfg:articulacion:canales},
\(\log q_+=-x-y\) and \(\log q_-=-x+y\);
adding and subtracting recovers \(x,y\).
Conversely, the recovered roots belong to \((0,1)\),
so that \(x\pm y=-\log q_\pm>0\); the resulting pair
remains in the chamber \(x>|y|\).
\end{proof}

In particular,
\(\widehat\mu\widehat\varepsilon=\widehat c^{-2}\) and
\(\widehat\mu/\widehat\varepsilon=\widehat Z^2\).
Exchanging the labels permutes \(q_+\) and \(q_-\),
preserves \(x\) and changes the sign of \(y\). Recovery
requires that orientation: the isolated product \(r_+r_-\)
preserves common propagation, but does not determine the two labeled
responses. The positive roots and cubic inversion recover
the channels within the fixed reader; the complete record of their
generation remains in the enriched state.

\subsection{The Catalan moment and functional recovery}
\label{hmtfg:articulacion:catalan}

The same incidences determine a relation between functions,
before either channel is evaluated. For \(0<s<1\), let
\(\mathcal D=s\,d/ds\), and define the oriented moment
\begin{equation}
 \mathcal C_2(s)
 =\int_0^s\log\frac{s}{t}\,
     \frac{1+t}{1+t+t^2}f(t)\,dt
 =\int_0^s\frac{\log(s/t)}{1+t^2}\,dt.
 \label{hmtfg:articulacion:integral}
\end{equation}
The second equality follows from the factorization
\(1+t+t^2+t^3=(1+t)(1+t^2)\). The continuous positive weight
\(dt/(1+t^2)\) is thus produced by the rational composition
of the constitutive reader, not by an identification with
\(\nu_0\).

\begin{proposition}[Differential inversion, uniqueness and endpoint reading]
\label{hmtfg:articulacion:restitucion-catalan}
The moment \(\mathcal C_2\) is the unique function
\(F\in C^2((0,1))\) satisfying
\begin{equation}
 \mathcal D^2F(s)=\frac{s(1+s)}{1+s+s^2}f(s),\qquad
 \lim_{s\downarrow0}F(s)=0.
 \label{hmtfg:articulacion:problema-diferencial}
\end{equation}
The constitutive response is recovered through
\begin{equation}
 \mathcal D^2\mathcal C_2(s)=\frac{s}{1+s^2},\qquad
 f(s)=\frac{1+s+s^2}{s(1+s)}\,
          \mathcal D^2\mathcal C_2(s).
 \label{hmtfg:articulacion:inversa-diferencial}
\end{equation}
Its continuous extension to the closed interval has the series
\begin{equation}
 \mathcal C_2(s)=\sum_{j=0}^{\infty}
       \frac{(-1)^j s^{2j+1}}{(2j+1)^2},\qquad 0\le s\le1,
 \qquad
 \mathcal C_2(1)=\sum_{j=0}^{\infty}\frac{(-1)^j}{(2j+1)^2}
 =G_{\rm Cat}.
 \label{hmtfg:articulacion:serie-catalan}
\end{equation}
\end{proposition}
\begin{proof}
The integrand of \eqref{hmtfg:articulacion:integral} is nonnegative
and
\[
 0\le\mathcal C_2(s)\le\int_0^s\log(s/t)\,dt=s.
\]
This proves integrability at zero and the required limit.
The identity \(\log(s/t)=\int_t^s du/u\), followed by
interchanging nonnegative integrals, provides
\[
 \mathcal C_2(s)
 =\int_0^s\frac1u\left(\int_0^u\frac{dt}{1+t^2}\right)du.
\]
On every compact interval of the open domain, the
fundamental theorem of calculus applies. Therefore,
\[
 \mathcal D\mathcal C_2(s)=\int_0^s\frac{dt}{1+t^2},
 \qquad
 \mathcal D^2\mathcal C_2(s)=\frac{s}{1+s^2}.
\]
The rational factorization used in
\eqref{hmtfg:articulacion:integral} gives
\eqref{hmtfg:articulacion:problema-diferencial}; solving for \(f\)
gives \eqref{hmtfg:articulacion:inversa-diferencial}, whose
denominators are positive in \((0,1)\).

If two solutions have the required zero limit, their difference
\(H\) satisfies \(\mathcal D^2H=0\). The substitution \(z=\log s\)
transforms this equation into \(d^2H(e^z)/dz^2=0\), whence
\(H(s)=a\log s+b\). Existence of the zero limit forces
\(a=b=0\), proving uniqueness.

For \(s<1\), the geometric expansion of \((1+t^2)^{-1}\)
can be integrated term by term with the logarithmic weight, since
\[
 \int_0^s t^{2j}\log(s/t)\,dt
 =\frac{s^{2j+1}}{(2j+1)^2},\qquad
 \sum_{j\ge0}\frac{s^{2j+1}}{(2j+1)^2}<\infty.
\]
This proves the series in the interior. The bound
\(1/(2j+1)^2\) gives absolute and uniform convergence of that
series on \([0,1]\). In the integral, the integrand extended
by zero for \(t>s\) is dominated by \(\log(1/t)\)
on \((0,1)\). Dominated convergence allows us to reach
\(s=1\) and agrees with the extension of the series. Its endpoint
value is precisely the stated Catalan series.
\end{proof}

Differential inversion uses the family \(\mathcal C_2(s)\),
not only its endpoint value \(G_{\rm Cat}\). In particular,
\(\mathcal D^2\mathcal C_2(s)\to1/2\) when \(s\uparrow1\).
The series differentiated twice with \(\mathcal D\) has
a radial Abel limit there; the series of the moment itself is
absolutely convergent at the boundary. This distinction preserves
the precise conditions of each passage to the limit.

\subsection{Composition of the recoveries and joint scope}
\label{hmtfg:articulacion:conclusion}

The preceding interrelations have explicit operations and
domains. The moments of the atomic measure \(\nu_0\)
produce \(J_{12}\), and its functional calculus recovers the profile
\(u_0\). The profile and first-root selector determine
the sequence whose scaling yields the readout \(\delta_{\mathrm F}\).
On the other hand, the constitutive function \(f\) produces
\(\mathcal C_2\) through the integral moment, and
\eqref{hmtfg:articulacion:inversa-diferencial} recovers \(f\)
from that family. Its labeled evaluations at
\(s_\pm=q_\pm^{30}\) determine the constitutive responses;
cubic inversion recovers both channels and their angular
coordinates. The endpoint reading of the moment gives \(G_{\rm Cat}\).

The measure \(\nu_0\) is positive, atomic and of rank twelve.
The Catalan moment integrates a continuous positive weight and has
an alternating expansion. The signs of that expansion are
coefficients of the geometric series; they are not negative weights of
\(\nu_0\), nor do they turn the positive integral into the same
spectral representation. Each reader retains its support,
orientation, normalization and evaluation operation.

Unity lies in the APP--TRIT--TPK structure that generates and
transports their antecedents; the functional relations are
established through the proved compositions. Equality of
origin does not identify operators with different domains.
In particular, this articulation does not postulate a scalar formula
determining \(\delta_{\mathrm F}\) from Catalan’s constant and
other physical constants. It places critical scaling alongside
effective recoveries between readers on the same support,
preserving the information required by each direction.
